\section{Minimal flows}\label{sec:minimal}

Minimal flows are the irreducible objects of topological dynamics: they cannot
be shrunk to a smaller subflow.

\begin{definition}[Minimal flow]\label{def:minimal}
\uses{def:flow}
A $G$-flow $X$ is \emph{minimal} if its only subflow is $X$ itself;
equivalently, the orbit $Gx$ is dense in $X$ for every $x\in X$.
\end{definition}

\begin{proposition}\label{prop:minimal-subflow}
\uses{def:minimal}
Every $G$-flow contains a minimal subflow.
\end{proposition}

\begin{proof}
Order the subflows of $X$ by reverse inclusion. The intersection of a chain of
subflows is closed and invariant, and it is nonempty by compactness, since the
chain has the finite intersection property. By Zorn's lemma there is a
minimal subflow.
\end{proof}

\begin{example}\label{ex:minimal}
\uses{def:minimal}
A fixed point is a minimal subflow. The rotation of the circle by an
irrational multiple of $2\pi$ is a minimal $\mathbb{Z}$-flow, since every orbit
is dense.
\end{example}

\begin{figure}[htbp]
\centering
\includegraphics[width=\linewidth]{irrational-rotation}
\caption{Example~\ref{ex:minimal}: the first $12$ and the first $100$ points $x$, $Rx$, $R^2x$, \dots{} of an orbit of the rotation $R$ of the circle by $2\pi\alpha$, where $\alpha=(\sqrt5-1)/2$ is irrational. Every orbit is dense, so the rotation is a minimal $\mathbb{Z}$-flow.}\label{fig:irrational-rotation}
\illustrates{ex:minimal}
\end{figure}
